\section{Related Work}

\subsection{Automatic Speech Quality Assessment}
Research on automated speech quality assessment has evolved from early CNN- and RNN-based predictors such as MOSNet and Quality-Net \citep{lo2019mosnet, fu2018quality} to self-supervised learning approaches based on wav2vec 2.0 and HuBERT \citep{baevski2020wav2vec, hsu2021hubert}, and more recently to Audio Language Models (ALMs) \citep{wang2025enabling, deshmukh2024pam, chen2025audio, zezario2025study}. Recent state-of-the-art MOS prediction systems, including UTMOS \citep{saeki2022utmos} and LE-SSL-MOS \citep{qi2023ssl}, have shown strong results on in-domain and OOD tracks of the VoiceMOS challenges. To address the inconsistencies in subjective ratings across datasets, recent work has explored bias-aware training losses \citep{mittag2021bias} and dataset score alignment frameworks such as AlignNet \citep{pieper2024alignnet}. While these methods improve cross-dataset performance, progress remains limited by heterogeneous MOS annotation protocols, motivating the need for a unified benchmark for consistent evaluation.

\subsection{Advances in Reward Modeling}
Recent advances in reward modeling (RM), originally introduced to align model outputs with human preferences \citep{ouyang2022training}, have led to diverse paradigms in text and vision. In the text domain, research has progressed from scalar reward models to generative reward modeling. The Critique-out-Loud (CLoud) framework \citep{ankner2024critique} introduces natural language critiques prior to scalar scoring, bridging generative judgment and reward modeling. More recently, \citet{liu2025inference} proposes Self-Principled Critique Tuning (SPCT), enabling GRMs to generate adaptive principles and critiques during inference, achieving state-of-the-art results across RM benchmarks. In the vision domain, RM has been extended to multimodal evaluation, with proprietary systems such as GPT-4V \citep{achiam2023gpt} demonstrating strong agreement with human judgments and open-source efforts like LLaVA-Critic \citep{xiong2025llava} unifying pointwise and pairwise scoring. By contrast, research on speech RM remains sparse, with no established frameworks for scalable or fine-grained preference alignment.